\documentclass[11pt]{article}

\usepackage[utf8]{inputenc}
\usepackage[T1]{fontenc}
\usepackage{lmodern}
\usepackage{geometry}
\usepackage{amsmath,amssymb,amsfonts,amsthm,mathtools}
\usepackage{bm}
\usepackage{enumitem}
\usepackage{microtype}
\usepackage{hyperref}
\usepackage{titlesec}
\usepackage{booktabs}
\usepackage{array}
\usepackage{setspace}
\usepackage{mathrsfs}
\usepackage{physics}

\geometry{margin=1in}
\hypersetup{colorlinks=true, linkcolor=black, urlcolor=blue, citecolor=black}
\setlist[enumerate]{topsep=0.4em,itemsep=0.35em}
\setlist[itemize]{topsep=0.3em,itemsep=0.25em}
\titleformat{\section}{\large\bfseries}{\thesection.}{0.5em}{}
\titleformat{\subsection}{\normalsize\bfseries}{\thesubsection.}{0.5em}{}
\setstretch{1.06}

\newcommand{\Aether}{\AE{}ther}
\newcommand{\Aflow}{\AE{}ther-flow}
\newcommand{\TheoryName}{\Aflow{} Interpretation of Relativity}
\newcommand{\TheoryProgram}{\Aether{} / \Aflow{} framework}

\newcommand{\CompletedMetric}{g^{\sharp}}
\newcommand{\CompletionCore}{\mathfrak C^{\mathrm{zr,aff}}}
\newcommand{\TraceCore}{\mathfrak C^{\mathrm{tr}}}
\newcommand{\AffineNucleus}{\mathfrak N^{\mathrm{aff}}}

\newcommand{\Car}{\mathrm{car}}
\newcommand{\Cur}{\mathrm{cur}}
\newcommand{\Hom}{\mathrm{hom}}

\newcommand{\ForSpace}[1]{\mathcal I_{#1}^{\mathrm{for}}}
\newcommand{\PrimShell}{\mathcal S_{\mathrm{prim}}}
\newcommand{\Reduction}[1]{\mathcal R_{#1}}
\newcommand{\CurrentCarrier}[1]{\mathfrak C_{#1}^{\mathrm{cur}}}
\newcommand{\EqCarrier}[1]{\mathfrak C_{#1}^{\mathrm{eq}}}
\newcommand{\MetricOut}{\mathscr G_{\mathrm{out}}}
\newcommand{\MatterOut}{\mathscr T_{\mathrm{out}}}

\newcommand{\VisibleAffine}[1]{\widehat X_{#1}}
\newcommand{\DataSpace}[1]{\mathcal Y_{#1}^{\mathrm{obs}}}
\newcommand{\AmbientTarget}[1]{E_{#1}^{\mathrm{amb}}}

\newcommand{\CoreState}[1]{\mathcal C_{#1}}
\newcommand{\CoreSpace}[1]{\mathcal Q_{#1}}
\newcommand{\CoreAffine}[1]{\mathcal F_{#1}^{\mathrm{co}}}
\newcommand{\CoreBody}[1]{Q_{#1}^{\mathrm{co}}}
\newcommand{\CoreStateProj}[1]{\Pi_{#1}^{\mathrm{co,st}}}
\newcommand{\CoreDataProj}[1]{\Pi_{#1}^{\mathrm{co,dat}}}
\newcommand{\CoreShellProj}[1]{\Pi_{#1}^{\mathrm{co,sh}}}
\newcommand{\CoreSym}[1]{\mathbb K_{#1}^{\mathrm{co}}}
\newcommand{\CoreTime}[1]{\tau_{#1}^{\mathrm{co}}}
\newcommand{\CoreEmbed}[1]{\eta_{#1}}

\newcommand{\GapSpace}[1]{H_{#1}}
\newcommand{\GapBody}[1]{D_{#1}}
\newcommand{\HiddenOperator}[1]{K_{#1}^{\mathrm{hid}}}

\newcommand{\TraceBody}[1]{P_{#1}^{\mathrm{tr}}}
\newcommand{\TraceDataMap}[1]{\Omega_{#1}^{\mathrm{tr,dat}}}
\newcommand{\TraceShellMap}[1]{\Omega_{#1}^{\mathrm{tr,sh}}}
\newcommand{\TraceSym}[1]{\mathbb K_{#1}^{\mathrm{tr}}}
\newcommand{\TraceTime}[1]{\tau_{#1}^{\mathrm{tr}}}
\newcommand{\TraceEmbed}[1]{\zeta_{#1}}

\newcommand{\RecAffine}[1]{\mathcal F_{#1}^{\mathrm{rec}}}
\newcommand{\RecBody}[1]{Q_{#1}^{\mathrm{rec}}}
\newcommand{\RecEmbed}[1]{\eta_{#1}^{\mathrm{rec}}}

\newcommand{\TraceShellLin}[1]{L_{#1}^{\mathrm{tr,sh}}}
\newcommand{\ShellKer}[1]{V_{#1}^{\mathrm{ker}}}
\newcommand{\ShellNormal}[1]{E_{#1}^{\mathrm{sh}}}
\newcommand{\ShellOffset}[1]{\lambda_{#1}^{\mathrm{sh}}}
\newcommand{\DirProj}[1]{\Pi_{#1}^{\parallel}}
\newcommand{\CoordSpace}[1]{\mathcal K_{#1}}
\newcommand{\CoordBody}[1]{P_{#1}}
\newcommand{\CoordIso}[1]{T_{#1}}
\newcommand{\AffineData}[1]{\widehat D_{#1}}
\newcommand{\ShellAffine}[1]{\mathcal A_{#1}^{\mathrm{sh}}}
\newcommand{\ShellEmbed}[1]{\chi_{#1}}
\newcommand{\SymGroup}[1]{\mathbb G_{#1}}
\newcommand{\TimeRev}[1]{\vartheta_{#1}}

\newtheorem{definition}{Definition}
\newtheorem{proposition}{Proposition}
\newtheorem{remark}{Remark}
\newtheorem{corollary}{Corollary}
\newtheorem{theorem}{Theorem}

\title{\textbf{The \TheoryName{}}\\[0.4em]
\large Primitive-Reservoir Observer-Reduction\\
\large Minimal Zero-Response Affine Completion Core\\
\large Collapse to a Minimal Completion-State Shell-Trace Core}
\author{}
\date{}

\begin{document}

\maketitle

\begin{abstract}
The current active theorem on the primitive-reservoir observer-reduction line already removed the exact-shell affine nucleus \(\AffineNucleus\) as the live unexplained stopping point: the nucleus is forced by the explicit minimal zero-response affine completion core \(\CompletionCore\). The next admissible benchmark-facing move recorded by the routing layer was therefore either to force \(\CompletionCore\) from still more primitive explicit substrate structure, or to prove a genuinely smaller collapse strictly inside \(\CompletionCore\).

The present note takes the second route. Rather than reopen a preserved deeper ambient-style relay and reargue its benchmark status, it proves that the full completion-core graph presentation carries no extra benchmark-facing freedom beyond a smaller completion-state shell-trace core \(\TraceCore\). For each sector, that smaller datum keeps only the completion-state body, the affine observer-data map, the affine shell-response map, the hidden body with positive gap operator, the induced symmetry and time-reversal structure, the exact zero-response shell realization in completion-state coordinates, fixed-data solvability, and the data-kernel coercive control. The affine completion plane and the completion body are then recovered canonically as the graph of those same trace maps over that same state body.

This is a genuine smaller theorem-level collapse strictly inside the current completion-core frontier. The shell-coordinate construction yielding the same exact-shell affine nucleus now factors through \(\TraceCore\), so the remaining honest burden is no longer the full completion-core graph presentation as such. The benchmark boundary remains fixed throughout: the one operative metric is still exactly \(\CompletedMetric\), the ordinary matter route is unchanged, there is no second operative metric, the low-energy matter law is not enlarged, and stronger first-principles promotion language remains paused.
\end{abstract}

\section{Introduction}

The active derivational program of \TheoryName{} has reached another point where depth alone does not count as progress. The direct completion-core forcing theorem already discharged the honest shell-facing move that was still open at the previous frontier: the exact-shell affine nucleus \(\AffineNucleus\) is a canonical output of the explicit minimal zero-response affine completion core \(\CompletionCore\). Once that is on record, another note that merely restates the nucleus or re-expands the larger primitive reversible constrained dynamics package no longer changes benchmark-facing status.

The current routing criterion left two admissible routes below that frontier. The first would be a theorem that forces \(\CompletionCore\) from still more primitive explicit substrate structure in a way that genuinely changes benchmark-facing status. The second would be a smaller theorem-level collapse strictly inside \(\CompletionCore\). The present note takes the second route explicitly. It does not promote an older deeper relay back to live-frontier status. Instead it asks a more precise internal question: inside the current completion-core package, what is genuinely independent benchmark-facing data, and what is merely affine-graph presentation once the completion-state trace has been fixed?

The answer is that the full completion plane and completion body are not independently live at current scope. They are reconstructed from a smaller completion-state shell-trace core consisting only of the completion-state body, the descended affine observer-data and shell-response maps, the hidden body with positive gap operator, the induced symmetry / time-reversal structure, the exact zero-response shell realization in completion-state coordinates, fixed-data solvability, and the same coercive control on data-invisible variations. The benchmark-relevant shell-coordinate construction of the current nucleus theorem already depends only on those smaller data.

\paragraph{Framework and claim boundary.}
\TheoryProgram{} denotes the broader \Aether{} / \Aflow{} framework built on the claims that the \Aether{} is the underlying four-dimensional substrate of reality, the \Aflow{} is its intrinsic ordered motion, observed three-dimensional space is the local experiential slice of that deeper substrate, S-time is the experienced order of change arising from matter, light, and the \Aflow{}, and observed expansion is the three-dimensional appearance of deeper four-dimensional ordered motion. Within that framework, \TheoryName{} denotes the exact-closure theory in which the effective gravitational dynamics are adopted to be exactly Einsteinian with universal matter coupling. In the active sequence, \emph{adoption} means use of that established relativistic dynamics without claiming substrate derivation, while \emph{derivation} is reserved for a first-principles recovery from explicit substrate variables.

The present note stays strictly inside that boundary. It does not introduce a second operative metric, it does not enlarge the low-energy matter law, and it does not upgrade adoption to completed first-principles derivation. Its narrower benchmark-facing role is only to compress the current completion-core burden to a smaller internal datum.

\section{Fixed Primitive Continuation Data}

\begin{definition}[Recorded primitive continuation data]
\label{def:fixed}
Fix the following continuation data from the active primitive-reservoir observer-reduction line.
\begin{enumerate}
    \item For each sector label \(\sigma\in\{\Car,\Cur,\Hom\}\), a primitive forcing shell
    \[
    \ForSpace{\sigma},
    \]
    a visible reduction map
    \[
    \Reduction{\sigma}:\ForSpace{\sigma}\to X_{\sigma},
    \]
    and shell-level current and equilibrium carriers
    \[
    \CurrentCarrier{\sigma}:\ForSpace{\sigma}\to Z_{\sigma}^{\mathrm{cur}},
    \qquad
    \EqCarrier{\sigma}:\ForSpace{\sigma}\to Z_{\sigma}^{\mathrm{eq}}.
    \]
    \item The product shell
    \[
    \PrimShell
    \equiv
    \ForSpace{\Car}\times \ForSpace{\Cur}\times \ForSpace{\Hom}.
    \]
    \item The recorded metric and ordinary-matter outputs
    \[
    \MetricOut:\PrimShell\to\{\text{observer metrics}\},
    \qquad
    \MatterOut:\PrimShell\times\Psi\to\{\text{observer stress tensors}\},
    \]
    satisfying
    \begin{equation}
    \MetricOut(\mathbf w)=\CompletedMetric,
    \qquad
    \MatterOut(\mathbf w,\psi)
    =
    T_{\mu\nu}^{\mathrm{matter}}[\CompletedMetric,\psi]
    \label{eq:fixedoutputs}
    \end{equation}
    for every \(\mathbf w\in\PrimShell\) and every matter configuration \(\psi\in\Psi\).
\end{enumerate}
\end{definition}

\begin{remark}
Definition \ref{def:fixed} is not reopened here. The primitive shell data, the one operative metric, and the ordinary matter route remain fixed continuation data. The present note only sharpens the internal benchmark-facing burden below the current completion-core frontier.
\end{remark}

\section{Current Frontier Package: Minimal Zero-Response Affine Completion Core}

\begin{definition}[Minimal zero-response affine completion core with hidden-sector gap]
\label{def:core}
Fix \(\sigma\in\{\Car,\Cur,\Hom\}\). A \emph{minimal zero-response affine completion core with hidden-sector gap} for sector \(\sigma\) consists of the following data.
\begin{enumerate}
    \item A finite-dimensional Euclidean completion-state space
    \[
    \CoreState{\sigma},
    \]
    an observer-data space
    \[
    \DataSpace{\sigma}
    \equiv
    \VisibleAffine{\sigma}\times Z_{\sigma}^{\mathrm{cur}}\times Z_{\sigma}^{\mathrm{eq}},
    \]
    an ambient shell-response space
    \[
    \AmbientTarget{\sigma},
    \]
    a hidden space
    \[
    \GapSpace{\sigma},
    \]
    and the zero-hidden completion space
    \[
    \CoreSpace{\sigma}
    \equiv
    \CoreState{\sigma}\oplus\DataSpace{\sigma}\oplus\AmbientTarget{\sigma}.
    \]
    Let
    \[
    \CoreStateProj{\sigma}:\CoreSpace{\sigma}\to\CoreState{\sigma},
    \qquad
    \CoreDataProj{\sigma}:\CoreSpace{\sigma}\to\DataSpace{\sigma},
    \qquad
    \CoreShellProj{\sigma}:\CoreSpace{\sigma}\to\AmbientTarget{\sigma}
    \]
    denote the coordinate projections.
    \item An affine completion plane
    \[
    \CoreAffine{\sigma}\subset\CoreSpace{\sigma}
    \]
    such that
    \[
    \CoreStateProj{\sigma}\big|_{\CoreAffine{\sigma}}:
    \CoreAffine{\sigma}\to\CoreState{\sigma}
    \]
    is an affine isomorphism.
    \item A compact convex completion body
    \[
    \CoreBody{\sigma}\subset\CoreAffine{\sigma}
    \]
    with nonempty relative interior in \(\CoreAffine{\sigma}\).
    \item A compact convex hidden body
    \[
    \GapBody{\sigma}\subset\GapSpace{\sigma}
    \]
    with \(0\in\operatorname{int}\GapBody{\sigma}\), together with a positive self-adjoint hidden-sector operator
    \[
    \HiddenOperator{\sigma}:\GapSpace{\sigma}\to\GapSpace{\sigma}
    \]
    satisfying
    \begin{equation}
    \ev{\eta,\HiddenOperator{\sigma}\eta}
    \ge
    \kappa_{\sigma}^{\mathrm{hid}}\|\eta\|^{2}
    \qquad
    \text{for all }\eta\in\GapSpace{\sigma}
    \label{eq:coregap}
    \end{equation}
    for some \(\kappa_{\sigma}^{\mathrm{hid}}>0\).
    \item A compact symmetry group
    \[
    \CoreSym{\sigma}\subset
    O(\CoreState{\sigma})\times
    O(\DataSpace{\sigma})\times
    O(\AmbientTarget{\sigma})\times
    O(\GapSpace{\sigma})
    \]
    and an orthogonal involution
    \[
    \CoreTime{\sigma}\in
    O(\CoreState{\sigma})\times
    O(\DataSpace{\sigma})\times
    O(\AmbientTarget{\sigma})\times
    O(\GapSpace{\sigma})
    \]
    preserving \(\CoreAffine{\sigma}\), \(\CoreBody{\sigma}\), the zero-response slice
    \[
    \CoreState{\sigma}\oplus\DataSpace{\sigma}\oplus\{0\},
    \]
    the hidden body \(\GapBody{\sigma}\), and \(\HiddenOperator{\sigma}\).
    \item A smooth embedding
    \[
    \jmath_{\sigma}:X_{\sigma}\hookrightarrow \VisibleAffine{\sigma}
    \]
    and a smooth zero-response shell realization
    \[
    \CoreEmbed{\sigma}:\ForSpace{\sigma}\hookrightarrow\CoreSpace{\sigma}
    \]
    whose image is exactly
    \[
    \CoreBody{\sigma}\cap
    \bigl(\CoreShellProj{\sigma}\bigr)^{-1}(0),
    \]
    and such that
    \begin{equation}
    \CoreDataProj{\sigma}\bigl(\CoreEmbed{\sigma}(w)\bigr)
    =
    \bigl(
    \jmath_{\sigma}(\Reduction{\sigma}(w)),
    \CurrentCarrier{\sigma}(w),
    \EqCarrier{\sigma}(w)
    \bigr)
    \label{eq:corecompat}
    \end{equation}
    for every \(w\in\ForSpace{\sigma}\).
    \item Fixed-data solvability on the zero-response slice:
    \begin{equation}
    \CoreDataProj{\sigma}\bigl(\CoreBody{\sigma}\bigr)
    =
    \CoreDataProj{\sigma}
    \Bigl(
    \CoreBody{\sigma}\cap
    \bigl(\CoreShellProj{\sigma}\bigr)^{-1}(0)
    \Bigr).
    \label{eq:coresolvability}
    \end{equation}
    \item Core-kernel coercivity on the zero-response shell: there exists \(\kappa_{\sigma}^{\mathrm{co}}>0\) such that for every
    \[
    w\in\ForSpace{\sigma},
    \qquad
    \delta c\in T_{\CoreEmbed{\sigma}(w)}\CoreAffine{\sigma}
    \]
    satisfying
    \[
    \CoreDataProj{\sigma}(\delta c)=0,
    \]
    one has
    \begin{equation}
    \bigl\|
    \CoreShellProj{\sigma}(\delta c)
    \bigr\|^{2}
    \ge
    \kappa_{\sigma}^{\mathrm{co}}
    \bigl\|
    \CoreStateProj{\sigma}(\delta c)
    \bigr\|^{2}.
    \label{eq:corecoercive}
    \end{equation}
\end{enumerate}
\end{definition}

\begin{remark}
Definition \ref{def:core} is the current internal burden below the direct exact-shell-affine-nucleus forcing theorem. The question addressed here is whether the full affine completion-plane presentation is itself larger than necessary once one asks only for the benchmark-relevant completion-state trace.
\end{remark}

\section{Smaller Internal Datum: Minimal Completion-State Shell-Trace Core}

\begin{definition}[Completion-state shell-trace core]
\label{def:trace}
Fix \(\sigma\in\{\Car,\Cur,\Hom\}\). A \emph{minimal completion-state shell-trace core with hidden-sector gap} for sector \(\sigma\) consists of the following data.
\begin{enumerate}
    \item A finite-dimensional Euclidean completion-state space
    \[
    \CoreState{\sigma},
    \]
    the same observer-data space \(\DataSpace{\sigma}\), shell-response space \(\AmbientTarget{\sigma}\), and hidden space \(\GapSpace{\sigma}\).
    \item A compact convex completion-state body
    \[
    \TraceBody{\sigma}\subset\CoreState{\sigma}
    \]
    with nonempty interior in \(\CoreState{\sigma}\).
    \item A compact convex hidden body
    \[
    \GapBody{\sigma}\subset\GapSpace{\sigma}
    \]
    with \(0\in\operatorname{int}\GapBody{\sigma}\), together with a positive self-adjoint hidden-sector operator
    \[
    \HiddenOperator{\sigma}:\GapSpace{\sigma}\to\GapSpace{\sigma}
    \]
    satisfying
    \begin{equation}
    \ev{\eta,\HiddenOperator{\sigma}\eta}
    \ge
    \kappa_{\sigma}^{\mathrm{hid}}\|\eta\|^{2}
    \qquad
    \text{for all }\eta\in\GapSpace{\sigma}.
    \label{eq:tracegap}
    \end{equation}
    \item Affine observer-data and shell-response maps
    \[
    \TraceDataMap{\sigma}:\CoreState{\sigma}\to\DataSpace{\sigma},
    \qquad
    \TraceShellMap{\sigma}:\CoreState{\sigma}\to\AmbientTarget{\sigma}.
    \]
    \item A compact symmetry group
    \[
    \TraceSym{\sigma}\subset
    O(\CoreState{\sigma})\times
    O(\DataSpace{\sigma})\times
    O(\AmbientTarget{\sigma})\times
    O(\GapSpace{\sigma})
    \]
    and an orthogonal involution
    \[
    \TraceTime{\sigma}\in
    O(\CoreState{\sigma})\times
    O(\DataSpace{\sigma})\times
    O(\AmbientTarget{\sigma})\times
    O(\GapSpace{\sigma})
    \]
    preserving \(\TraceBody{\sigma}\), the exact zero-response trace slice
    \[
    \TraceBody{\sigma}\cap\bigl(\TraceShellMap{\sigma}\bigr)^{-1}(0),
    \]
    the hidden body \(\GapBody{\sigma}\), and \(\HiddenOperator{\sigma}\), and satisfying
    \[
    \TraceDataMap{\sigma}(g_{\mathrm{st}}m)
    =
    g_{\mathrm{dat}}\TraceDataMap{\sigma}(m),
    \qquad
    \TraceShellMap{\sigma}(g_{\mathrm{st}}m)
    =
    g_{\mathrm{sh}}\TraceShellMap{\sigma}(m)
    \]
    for every \(m\in\CoreState{\sigma}\) and every
    \[
    (g_{\mathrm{st}},g_{\mathrm{dat}},g_{\mathrm{sh}},g_{\mathrm{hid}})
    \in\TraceSym{\sigma},
    \]
    with the analogous intertwining equalities for the components of \(\TraceTime{\sigma}\).
    \item A smooth embedding
    \[
    \jmath_{\sigma}:X_{\sigma}\hookrightarrow\VisibleAffine{\sigma}
    \]
    and a smooth exact zero-response trace realization
    \[
    \TraceEmbed{\sigma}:\ForSpace{\sigma}\hookrightarrow\CoreState{\sigma}
    \]
    whose image is exactly
    \[
    \TraceBody{\sigma}\cap
    \bigl(\TraceShellMap{\sigma}\bigr)^{-1}(0),
    \]
    and such that
    \begin{equation}
    \TraceDataMap{\sigma}\bigl(\TraceEmbed{\sigma}(w)\bigr)
    =
    \bigl(
    \jmath_{\sigma}(\Reduction{\sigma}(w)),
    \CurrentCarrier{\sigma}(w),
    \EqCarrier{\sigma}(w)
    \bigr)
    \label{eq:tracecompat}
    \end{equation}
    for every \(w\in\ForSpace{\sigma}\).
    \item Fixed-data solvability on the trace body:
    \begin{equation}
    \TraceDataMap{\sigma}\bigl(\TraceBody{\sigma}\bigr)
    =
    \TraceDataMap{\sigma}
    \Bigl(
    \TraceBody{\sigma}\cap
    \bigl(\TraceShellMap{\sigma}\bigr)^{-1}(0)
    \Bigr).
    \label{eq:tracesolvability}
    \end{equation}
    \item Trace-kernel coercivity: there exists \(\kappa_{\sigma}^{\mathrm{tr}}>0\) such that for every
    \[
    \delta m\in\CoreState{\sigma}
    \]
    satisfying
    \[
    D\TraceDataMap{\sigma}\,\delta m=0,
    \]
    one has
    \begin{equation}
    \bigl\|
    D\TraceShellMap{\sigma}\,\delta m
    \bigr\|^{2}
    \ge
    \kappa_{\sigma}^{\mathrm{tr}}\|\delta m\|^{2}.
    \label{eq:tracecoercive}
    \end{equation}
\end{enumerate}
\end{definition}

\begin{remark}
Definition \ref{def:trace} is strictly smaller than Definition \ref{def:core} in the precise sense relevant here. It does not retain a separate affine completion plane or a separate completion body inside a larger graph space. It retains only the completion-state body together with the descended affine data and shell maps and the hidden gap package.
\end{remark}

\section{Canonical Trace Extraction from the Completion Core}

\begin{definition}[Canonical trace core extracted from a completion core]
\label{def:extract}
Assume Definition \ref{def:core}. Define the associated completion-state shell-trace core as follows.
\begin{enumerate}
    \item The completion-state body is
    \begin{equation}
    \TraceBody{\sigma}
    \equiv
    \CoreStateProj{\sigma}\bigl(\CoreBody{\sigma}\bigr).
    \label{eq:tracebody}
    \end{equation}
    \item The affine observer-data and shell-response trace maps are
    \begin{equation}
    \TraceDataMap{\sigma}
    \equiv
    \CoreDataProj{\sigma}\circ
    \bigl(
    \CoreStateProj{\sigma}\big|_{\CoreAffine{\sigma}}
    \bigr)^{-1},
    \qquad
    \TraceShellMap{\sigma}
    \equiv
    \CoreShellProj{\sigma}\circ
    \bigl(
    \CoreStateProj{\sigma}\big|_{\CoreAffine{\sigma}}
    \bigr)^{-1}.
    \label{eq:tracemaps}
    \end{equation}
    \item The symmetry group, time reversal, hidden body, and hidden operator are retained unchanged:
    \[
    \TraceSym{\sigma}\equiv\CoreSym{\sigma},
    \qquad
    \TraceTime{\sigma}\equiv\CoreTime{\sigma}.
    \]
    \item The exact zero-response trace realization is
    \begin{equation}
    \TraceEmbed{\sigma}
    \equiv
    \CoreStateProj{\sigma}\circ\CoreEmbed{\sigma}.
    \label{eq:traceembed}
    \end{equation}
\end{enumerate}
\end{definition}

\begin{proposition}[Every completion core canonically yields a trace core]
\label{prop:trace}
Assume Definitions \ref{def:core} and \ref{def:extract}. Then the extracted data satisfy every requirement of Definition \ref{def:trace}.
\end{proposition}

\begin{proof}
Because \(\CoreStateProj{\sigma}\big|_{\CoreAffine{\sigma}}\) is an affine isomorphism and \(\CoreBody{\sigma}\subset\CoreAffine{\sigma}\) is compact and convex with nonempty relative interior, the projected state body \(\TraceBody{\sigma}\) is compact and convex with nonempty interior in \(\CoreState{\sigma}\). The hidden body and hidden operator are unchanged, so \eqref{eq:tracegap} is exactly \eqref{eq:coregap}.

The maps \(\TraceDataMap{\sigma}\) and \(\TraceShellMap{\sigma}\) are affine by construction. If
\[
(g_{\mathrm{st}},g_{\mathrm{dat}},g_{\mathrm{sh}},g_{\mathrm{hid}})
\in\CoreSym{\sigma},
\qquad
m\in\CoreState{\sigma},
\]
then the point
\[
\bigl(
m,\TraceDataMap{\sigma}(m),\TraceShellMap{\sigma}(m)
\bigr)
\in\CoreAffine{\sigma}
\]
is carried by the symmetry to
\[
\bigl(
g_{\mathrm{st}}m,\,
g_{\mathrm{dat}}\TraceDataMap{\sigma}(m),\,
g_{\mathrm{sh}}\TraceShellMap{\sigma}(m)
\bigr)\in\CoreAffine{\sigma},
\]
so the defining graph property of \eqref{eq:tracemaps} yields the required intertwining equalities. The same argument applies to \(\TraceTime{\sigma}\). Because the completion-core symmetries preserve \(\CoreBody{\sigma}\), the shell-zero slice, \(\GapBody{\sigma}\), and \(\HiddenOperator{\sigma}\), the extracted trace symmetries preserve the corresponding trace-core data as well.

For the exact zero-response trace realization, the image of \(\TraceEmbed{\sigma}\) is contained in
\[
\TraceBody{\sigma}\cap
\bigl(\TraceShellMap{\sigma}\bigr)^{-1}(0)
\]
because \(\CoreEmbed{\sigma}(\ForSpace{\sigma})\subset\CoreBody{\sigma}\cap(\CoreShellProj{\sigma})^{-1}(0)\). Conversely, if
\[
m\in\TraceBody{\sigma}\cap
\bigl(\TraceShellMap{\sigma}\bigr)^{-1}(0),
\]
then
\[
c\equiv
\bigl(
\CoreStateProj{\sigma}\big|_{\CoreAffine{\sigma}}
\bigr)^{-1}(m)
\in\CoreBody{\sigma}
\]
and \(\CoreShellProj{\sigma}(c)=0\). Hence \(c\in\CoreBody{\sigma}\cap(\CoreShellProj{\sigma})^{-1}(0)\), so by Definition \ref{def:core} there is a unique \(w\in\ForSpace{\sigma}\) with \(c=\CoreEmbed{\sigma}(w)\). Therefore \(m=\TraceEmbed{\sigma}(w)\). Equation \eqref{eq:tracecompat} follows directly from \eqref{eq:corecompat}.

For trace solvability, let \(u\in\TraceDataMap{\sigma}(\TraceBody{\sigma})\). Choose \(c\in\CoreBody{\sigma}\) with
\[
u=\CoreDataProj{\sigma}(c).
\]
By \eqref{eq:coresolvability}, there exists
\[
c_{0}\in\CoreBody{\sigma}\cap
\bigl(\CoreShellProj{\sigma}\bigr)^{-1}(0)
\]
with \(\CoreDataProj{\sigma}(c_{0})=u\). Set
\[
m_{0}\equiv\CoreStateProj{\sigma}(c_{0}).
\]
Then \(m_{0}\in\TraceBody{\sigma}\cap(\TraceShellMap{\sigma})^{-1}(0)\) and \(\TraceDataMap{\sigma}(m_{0})=u\), proving \eqref{eq:tracesolvability}.

Finally, let \(\delta m\in\CoreState{\sigma}\) satisfy \(D\TraceDataMap{\sigma}\,\delta m=0\). Define the corresponding tangent vector to the completion graph by
\[
\delta c
\equiv
\bigl(
\delta m,\,
D\TraceDataMap{\sigma}\,\delta m,\,
D\TraceShellMap{\sigma}\,\delta m
\bigr)
\in T\CoreAffine{\sigma}.
\]
Because \(\CoreAffine{\sigma}\) is affine, this tangent space is the same at every shell point, so \(\delta c\in T_{\CoreEmbed{\sigma}(w)}\CoreAffine{\sigma}\) for any \(w\in\ForSpace{\sigma}\). The data-kernel assumption gives \(\CoreDataProj{\sigma}(\delta c)=0\). Applying \eqref{eq:corecoercive} yields
\[
\bigl\|D\TraceShellMap{\sigma}\,\delta m\bigr\|^{2}
=
\bigl\|\CoreShellProj{\sigma}(\delta c)\bigr\|^{2}
\ge
\kappa_{\sigma}^{\mathrm{co}}
\bigl\|\CoreStateProj{\sigma}(\delta c)\bigr\|^{2}
=
\kappa_{\sigma}^{\mathrm{co}}\|\delta m\|^{2}.
\]
This is \eqref{eq:tracecoercive} with \(\kappa_{\sigma}^{\mathrm{tr}}=\kappa_{\sigma}^{\mathrm{co}}\).
\end{proof}

\section{Reconstruction of the Full Completion Core from the Trace Core}

\begin{definition}[Canonical completion-core reconstruction from a trace core]
\label{def:reconstruct}
Assume Definition \ref{def:trace}. Define the reconstructed completion-core data by
\begin{align}
\CoreSpace{\sigma}
&\equiv
\CoreState{\sigma}\oplus\DataSpace{\sigma}\oplus\AmbientTarget{\sigma},
\nonumber\\
\RecAffine{\sigma}
&\equiv
\left\{
\bigl(
m,\TraceDataMap{\sigma}(m),\TraceShellMap{\sigma}(m)
\bigr):
m\in\CoreState{\sigma}
\right\},
\label{eq:recaffine}\\
\RecBody{\sigma}
&\equiv
\left\{
\bigl(
m,\TraceDataMap{\sigma}(m),\TraceShellMap{\sigma}(m)
\bigr):
m\in\TraceBody{\sigma}
\right\},
\label{eq:recbody}
\end{align}
and
\begin{equation}
\RecEmbed{\sigma}(w)
\equiv
\bigl(
\TraceEmbed{\sigma}(w),
\TraceDataMap{\sigma}(\TraceEmbed{\sigma}(w)),
0
\bigr).
\label{eq:recembed}
\end{equation}
Use the standard coordinate projections on \(\CoreSpace{\sigma}\), and keep the same hidden body, hidden operator, symmetry group, and time reversal.
\end{definition}

\begin{proposition}[A trace core reconstructs a completion core]
\label{prop:reconstruct}
Assume Definitions \ref{def:trace} and \ref{def:reconstruct}. Then the reconstructed data satisfy every requirement of Definition \ref{def:core}.
\end{proposition}

\begin{proof}
By \eqref{eq:recaffine}, the reconstructed affine plane is the graph of the affine pair \((\TraceDataMap{\sigma},\TraceShellMap{\sigma})\), so the state projection is an affine isomorphism from \(\RecAffine{\sigma}\) onto \(\CoreState{\sigma}\). Because \(\TraceBody{\sigma}\) is compact and convex with nonempty interior, its graph image \(\RecBody{\sigma}\) is compact and convex with nonempty relative interior in \(\RecAffine{\sigma}\). The hidden body and hidden operator already satisfy \eqref{eq:coregap} by \eqref{eq:tracegap}.

The symmetry and time-reversal data preserve the reconstructed graph because the trace-core intertwining equalities imply
\[
\bigl(
g_{\mathrm{st}}m,\,
g_{\mathrm{dat}}\TraceDataMap{\sigma}(m),\,
g_{\mathrm{sh}}\TraceShellMap{\sigma}(m)
\bigr)
\in\RecAffine{\sigma}
\]
whenever \(m\in\CoreState{\sigma}\) and
\[
(g_{\mathrm{st}},g_{\mathrm{dat}},g_{\mathrm{sh}},g_{\mathrm{hid}})
\in\TraceSym{\sigma}.
\]
They preserve \(\RecBody{\sigma}\), the zero-response slice, \(\GapBody{\sigma}\), and \(\HiddenOperator{\sigma}\) for the same reason.

For the shell realization, \(\RecEmbed{\sigma}(\ForSpace{\sigma})\) is contained in
\[
\RecBody{\sigma}\cap
\bigl(\CoreShellProj{\sigma}\bigr)^{-1}(0)
\]
because \(\TraceEmbed{\sigma}(\ForSpace{\sigma})\subset\TraceBody{\sigma}\cap(\TraceShellMap{\sigma})^{-1}(0)\). Conversely, if
\[
(m,u,0)\in
\RecBody{\sigma}\cap
\bigl(\CoreShellProj{\sigma}\bigr)^{-1}(0),
\]
then \(m\in\TraceBody{\sigma}\), \(u=\TraceDataMap{\sigma}(m)\), and \(\TraceShellMap{\sigma}(m)=0\). Hence \(m=\TraceEmbed{\sigma}(w)\) for a unique \(w\in\ForSpace{\sigma}\), and therefore
\[
(m,u,0)=\RecEmbed{\sigma}(w).
\]
Equation \eqref{eq:corecompat} follows from \eqref{eq:tracecompat}.

For fixed-data solvability, let \(u\in\CoreDataProj{\sigma}(\RecBody{\sigma})\). Then \(u=\TraceDataMap{\sigma}(m)\) for some \(m\in\TraceBody{\sigma}\). By \eqref{eq:tracesolvability}, there exists
\[
m_{0}\in\TraceBody{\sigma}\cap
\bigl(\TraceShellMap{\sigma}\bigr)^{-1}(0)
\]
with \(\TraceDataMap{\sigma}(m_{0})=u\). Then
\[
\bigl(
m_{0},u,0
\bigr)
\in
\RecBody{\sigma}\cap
\bigl(\CoreShellProj{\sigma}\bigr)^{-1}(0),
\]
which proves \eqref{eq:coresolvability}.

Finally, let
\[
\delta c\in T_{\RecEmbed{\sigma}(w)}\RecAffine{\sigma}
\]
with \(\CoreDataProj{\sigma}(\delta c)=0\). Because \(\RecAffine{\sigma}\) is the graph of the affine trace maps, there exists a unique \(\delta m\in\CoreState{\sigma}\) such that
\[
\delta c=
\bigl(
\delta m,\,
D\TraceDataMap{\sigma}\,\delta m,\,
D\TraceShellMap{\sigma}\,\delta m
\bigr).
\]
The data-kernel condition is \(D\TraceDataMap{\sigma}\,\delta m=0\). Hence \eqref{eq:tracecoercive} gives
\[
\bigl\|
\CoreShellProj{\sigma}(\delta c)
\bigr\|^{2}
=
\bigl\|
D\TraceShellMap{\sigma}\,\delta m
\bigr\|^{2}
\ge
\kappa_{\sigma}^{\mathrm{tr}}\|\delta m\|^{2}
=
\kappa_{\sigma}^{\mathrm{tr}}
\bigl\|
\CoreStateProj{\sigma}(\delta c)
\bigr\|^{2}.
\]
This is \eqref{eq:corecoercive}.
\end{proof}

\begin{proposition}[The extracted trace core reconstructs the original completion core exactly]
\label{prop:exact}
Assume Definitions \ref{def:core}, \ref{def:extract}, and \ref{def:reconstruct}. Then the reconstructed completion-core data are exactly the original ones:
\[
\RecAffine{\sigma}=\CoreAffine{\sigma},
\qquad
\RecBody{\sigma}=\CoreBody{\sigma},
\qquad
\RecEmbed{\sigma}=\CoreEmbed{\sigma}.
\]
\end{proposition}

\begin{proof}
By definition of \(\TraceDataMap{\sigma}\) and \(\TraceShellMap{\sigma}\), every point \(c\in\CoreAffine{\sigma}\) with
\[
m=\CoreStateProj{\sigma}(c)
\]
satisfies
\[
c=
\bigl(
m,\TraceDataMap{\sigma}(m),\TraceShellMap{\sigma}(m)
\bigr),
\]
so \(c\in\RecAffine{\sigma}\). Conversely, every point of \(\RecAffine{\sigma}\) is obtained by applying
\[
\bigl(
\CoreStateProj{\sigma}\big|_{\CoreAffine{\sigma}}
\bigr)^{-1}
\]
to some \(m\in\CoreState{\sigma}\), hence lies in \(\CoreAffine{\sigma}\). Therefore \(\RecAffine{\sigma}=\CoreAffine{\sigma}\).

If \(c\in\CoreBody{\sigma}\), then with \(m=\CoreStateProj{\sigma}(c)\in\TraceBody{\sigma}\) the previous identity shows \(c\in\RecBody{\sigma}\). Conversely, let
\[
\bigl(
m,\TraceDataMap{\sigma}(m),\TraceShellMap{\sigma}(m)
\bigr)\in\RecBody{\sigma}.
\]
Because \(m\in\TraceBody{\sigma}=\CoreStateProj{\sigma}(\CoreBody{\sigma})\), there exists \(c\in\CoreBody{\sigma}\) with \(\CoreStateProj{\sigma}(c)=m\). Since \(c\in\CoreAffine{\sigma}\) and \(\CoreAffine{\sigma}\) is exactly the graph of the extracted trace maps, \(c\) must equal the displayed graph point. Hence \(\RecBody{\sigma}=\CoreBody{\sigma}\).

Finally,
\[
\RecEmbed{\sigma}(w)
=
\bigl(
\CoreStateProj{\sigma}(\CoreEmbed{\sigma}(w)),
\CoreDataProj{\sigma}(\CoreEmbed{\sigma}(w)),
0
\bigr)
=
\CoreEmbed{\sigma}(w)
\]
because \(\CoreEmbed{\sigma}(w)\) already lies in the shell-zero slice.
\end{proof}

\section{Downstream Exact-Shell Affine Structure Factors Through the Trace Core}

\begin{definition}[Canonical shell coordinates from the trace core]
\label{def:shell}
Assume Definition \ref{def:trace}. Define the shell coordinates induced by the trace core as follows.
\begin{enumerate}
    \item Let
    \[
    \TraceShellLin{\sigma}:\CoreState{\sigma}\to\AmbientTarget{\sigma}
    \]
    denote the linear part of the affine map \(\TraceShellMap{\sigma}\). Define
    \[
    \ShellKer{\sigma}\equiv\ker\TraceShellLin{\sigma},
    \qquad
    \ShellNormal{\sigma}\equiv\operatorname{im}\TraceShellLin{\sigma}.
    \]
    Let
    \[
    \DirProj{\sigma}:\CoreState{\sigma}\to\ShellKer{\sigma}
    \]
    be the orthogonal projection onto \(\ShellKer{\sigma}\), and equip \(\ShellNormal{\sigma}\) with the Euclidean structure transported from \((\ShellKer{\sigma})^{\perp}\) by
    \[
    \TraceShellLin{\sigma}\big|_{(\ShellKer{\sigma})^{\perp}}
    :
    (\ShellKer{\sigma})^{\perp}\xrightarrow{\cong}\ShellNormal{\sigma}.
    \]
    \item Choose the shell offset \(\ShellOffset{\sigma}\in\ShellNormal{\sigma}\) by the affine decomposition
    \begin{equation}
    \TraceShellMap{\sigma}(m)
    =
    \TraceShellLin{\sigma}(m)-\ShellOffset{\sigma}.
    \label{eq:traceshelldecomp}
    \end{equation}
    \item The shell-coordinate space is
    \[
    \CoordSpace{\sigma}\equiv
    \ShellKer{\sigma}\oplus\ShellNormal{\sigma},
    \]
    with product Euclidean structure, and the shell-coordinate isomorphism is
    \begin{equation}
    \CoordIso{\sigma}(m)
    \equiv
    \bigl(
    \DirProj{\sigma}m,\,
    \TraceShellLin{\sigma}(m)
    \bigr).
    \label{eq:tracecoordiso}
    \end{equation}
    \item The shell-coordinate body and shell-data readout are
    \begin{equation}
    \CoordBody{\sigma}
    \equiv
    \CoordIso{\sigma}\bigl(\TraceBody{\sigma}\bigr),
    \qquad
    \AffineData{\sigma}
    \equiv
    \TraceDataMap{\sigma}\circ\CoordIso{\sigma}^{-1}.
    \label{eq:tracecoordbody}
    \end{equation}
    \item The exact shell affine plane is
    \begin{equation}
    \ShellAffine{\sigma}
    \equiv
    \left\{
    (v,e)\in\CoordSpace{\sigma}:
    e=\ShellOffset{\sigma}
    \right\},
    \label{eq:traceshellplane}
    \end{equation}
    and the exact shell realization is
    \begin{equation}
    \ShellEmbed{\sigma}(w)
    \equiv
    \CoordIso{\sigma}\bigl(\TraceEmbed{\sigma}(w)\bigr).
    \label{eq:traceshellembed}
    \end{equation}
    \item The induced symmetry group and time reversal are
    \[
    \SymGroup{\sigma}
    \equiv
    \left\{
    \bigl(
    (g_{\mathrm{st}}|_{\ShellKer{\sigma}},\,g_{\mathrm{sh}}|_{\ShellNormal{\sigma}}),
    g_{\mathrm{hid}}
    \bigr):
    (g_{\mathrm{st}},g_{\mathrm{dat}},g_{\mathrm{sh}},g_{\mathrm{hid}})
    \in\TraceSym{\sigma}
    \right\},
    \]
    \[
    \TimeRev{\sigma}
    \equiv
    \bigl(
    (\TraceTime{\sigma}|_{\ShellKer{\sigma}},
    \TraceTime{\sigma}|_{\ShellNormal{\sigma}}),
    \TraceTime{\sigma}|_{\GapSpace{\sigma}}
    \bigr).
    \]
\end{enumerate}
\end{definition}

\begin{proposition}[The current exact-shell affine nucleus is already carried by the trace core]
\label{prop:nucleus}
Assume Definitions \ref{def:trace} and \ref{def:shell}. Then the shell-coordinate body \(\CoordBody{\sigma}\), affine shell-data readout \(\AffineData{\sigma}\), exact shell affine plane \(\ShellAffine{\sigma}\), exact shell realization \(\ShellEmbed{\sigma}\), hidden body \(\GapBody{\sigma}\), hidden operator \(\HiddenOperator{\sigma}\), induced symmetry / time-reversal structure, fixed-data shell solvability, and affine-kernel coercive control are all determined by the trace core itself. In particular, the shell-coordinate construction yielding the current exact-shell affine nucleus \(\AffineNucleus\) factors through the trace core.
\end{proposition}

\begin{proof}
Because \(\TraceShellLin{\sigma}\big|_{(\ShellKer{\sigma})^{\perp}}\) is a linear isomorphism onto \(\ShellNormal{\sigma}\), the map \(\CoordIso{\sigma}\) is a linear isomorphism; with the chosen Euclidean structures it is an isometry. Hence \(\CoordBody{\sigma}\) is compact and convex with nonempty interior whenever \(\TraceBody{\sigma}\) is.

If \(w\in\ForSpace{\sigma}\), then \(\TraceEmbed{\sigma}(w)\in\TraceBody{\sigma}\) and \(\TraceShellMap{\sigma}(\TraceEmbed{\sigma}(w))=0\). By \eqref{eq:traceshelldecomp},
\[
\TraceShellLin{\sigma}(\TraceEmbed{\sigma}(w))
=
\ShellOffset{\sigma},
\]
so \(\ShellEmbed{\sigma}(w)\in\CoordBody{\sigma}\cap\ShellAffine{\sigma}\). Conversely, if \(q=\CoordIso{\sigma}(m)\in\CoordBody{\sigma}\cap\ShellAffine{\sigma}\), then \(m\in\TraceBody{\sigma}\) and the second coordinate condition forces \(\TraceShellLin{\sigma}(m)=\ShellOffset{\sigma}\), hence \(\TraceShellMap{\sigma}(m)=0\). Therefore \(m=\TraceEmbed{\sigma}(w)\) for a unique \(w\in\ForSpace{\sigma}\), and so \(q=\ShellEmbed{\sigma}(w)\). Thus \(\ShellEmbed{\sigma}\) is a diffeomorphism onto the full exact shell slice of \(\CoordBody{\sigma}\).

Equation \eqref{eq:tracecompat} immediately gives
\[
\AffineData{\sigma}\bigl(\ShellEmbed{\sigma}(w)\bigr)
=
\TraceDataMap{\sigma}\bigl(\TraceEmbed{\sigma}(w)\bigr)
=
\bigl(
\jmath_{\sigma}(\Reduction{\sigma}(w)),
\CurrentCarrier{\sigma}(w),
\EqCarrier{\sigma}(w)
\bigr).
\]
The hidden body, hidden operator, and induced symmetry / time-reversal structure are already part of the trace core.

For shell solvability, let \(u\in\AffineData{\sigma}(\CoordBody{\sigma})\). Then \(u=\TraceDataMap{\sigma}(m)\) for some \(m\in\TraceBody{\sigma}\). By \eqref{eq:tracesolvability}, there exists
\[
m_{0}\in\TraceBody{\sigma}\cap
\bigl(\TraceShellMap{\sigma}\bigr)^{-1}(0)
\]
with \(\TraceDataMap{\sigma}(m_{0})=u\). Hence
\[
\CoordIso{\sigma}(m_{0})\in\CoordBody{\sigma}\cap\ShellAffine{\sigma}
\]
and \(\AffineData{\sigma}(\CoordIso{\sigma}(m_{0}))=u\).

Finally, let
\[
\delta q\in\CoordSpace{\sigma}
\]
satisfy \(D\AffineData{\sigma}\,\delta q=0\), and let \(\delta m\equiv\CoordIso{\sigma}^{-1}(\delta q)\). Since \(\CoordIso{\sigma}\) is an isometry,
\[
\|\delta q\|=\|\delta m\|.
\]
Moreover,
\[
D\AffineData{\sigma}\,\delta q
=
D\TraceDataMap{\sigma}\,\delta m,
\]
so the trace-kernel condition gives
\[
\bigl\|
D\TraceShellMap{\sigma}\,\delta m
\bigr\|^{2}
\ge
\kappa_{\sigma}^{\mathrm{tr}}\|\delta m\|^{2}.
\]
But the second component of \(\delta q\) is exactly \(D\TraceShellMap{\sigma}\,\delta m\). Therefore the shell-normal part of \(\delta q\) obeys the same coercive lower bound, and adjoining any hidden variation \(\delta n\in\GapSpace{\sigma}\) gives the same affine-kernel coercive control as in the current nucleus theorem. Thus every downstream shell-coordinate constituent relevant to \(\AffineNucleus\) is already carried by the trace core.
\end{proof}

\begin{proposition}[Separate completion-plane presentation is benchmark neutral at current scope]
\label{prop:neutral}
Assume Definition \ref{def:core}. Then the separate presentation of the full affine completion plane \(\CoreAffine{\sigma}\), the full completion body \(\CoreBody{\sigma}\), and the ambient graph coordinates of the completion core adds no independent benchmark-facing freedom beyond the smaller trace core extracted in Definition \ref{def:extract}.
\end{proposition}

\begin{proof}
By Proposition \ref{prop:exact}, the full completion-core graph presentation is reconstructed exactly from the trace core. By Proposition \ref{prop:nucleus}, all downstream shell-coordinate data relevant to the current exact-shell affine nucleus are already determined by that same trace core. Therefore the separate completion-plane and completion-body presentation does not contribute new benchmark-facing information beyond the trace-core datum.
\end{proof}

\begin{proposition}[The benchmark package remains fixed]
\label{prop:benchmark}
Every completion-state shell-trace core obtained here preserves the same benchmark one-metric package \eqref{eq:fixedoutputs}. In particular, the operative metric remains exactly \(\CompletedMetric\), the ordinary matter route is unchanged, there is no second operative metric, and the low-energy matter law is not enlarged.
\end{proposition}

\begin{proof}
The present note removes only internal presentational redundancy inside the current completion-core frontier. It does not alter the primitive shell \(\PrimShell\), the recorded metric map \(\MetricOut\), or the recorded matter map \(\MatterOut\). Therefore \eqref{eq:fixedoutputs} remains exactly in force.
\end{proof}

\section{Main Theorem}

\begin{theorem}[Minimal zero-response affine completion core collapses to a minimal completion-state shell-trace core]
\label{thm:main}
Assume the fixed primitive continuation data of Definition \ref{def:fixed} and, for each sector \(\sigma\in\{\Car,\Cur,\Hom\}\), a minimal zero-response affine completion core in the sense of Definition \ref{def:core}. Then:
\begin{enumerate}
    \item the completion core canonically yields a smaller completion-state shell-trace core by Proposition \ref{prop:trace};
    \item the full affine completion-plane presentation is reconstructed exactly from that smaller trace core by Propositions \ref{prop:reconstruct} and \ref{prop:exact};
    \item the shell-coordinate construction yielding the current exact-shell affine nucleus \(\AffineNucleus\) already factors through that smaller trace core by Proposition \ref{prop:nucleus};
    \item the separate completion-plane and completion-body presentation therefore adds no independent benchmark-facing freedom beyond the trace core by Proposition \ref{prop:neutral};
    \item the benchmark one-metric package remains exactly the same by Proposition \ref{prop:benchmark};
    \item consequently the live internal burden below the direct completion-core frontier collapses from the full minimal zero-response affine completion core \(\CompletionCore\) to the smaller completion-state shell-trace core \(\TraceCore\). The remaining honest burden is therefore no longer the full completion-core graph presentation as such. What remains is to derive or uniquely force \(\TraceCore\) from still more primitive explicit substrate structure, or else prove a still smaller theorem-level collapse strictly inside \(\TraceCore\), while preserving the same benchmark package and claim boundary.
\end{enumerate}
\end{theorem}

\begin{proof}
Item (1) is Proposition \ref{prop:trace}. Item (2) is Propositions \ref{prop:reconstruct} and \ref{prop:exact}. Item (3) is Proposition \ref{prop:nucleus}. Item (4) is Proposition \ref{prop:neutral}. Item (5) is Proposition \ref{prop:benchmark}. Item (6) is the routing consequence of the previous items together with the current safeguard against counting another same-output deeper-origin relay as benchmark-facing progress when the benchmark-relevant internal datum has already been sharpened.
\end{proof}

\begin{corollary}[What must be done next]
\label{cor:next}
After Theorem \ref{thm:main}, the next honest benchmark-facing manuscript must either:
\begin{enumerate}
    \item derive or uniquely force the completion-state shell-trace core \(\TraceCore\) from still more primitive explicit substrate structure in a way that changes benchmark-facing status;
    \item identify a genuinely smaller theorem-level datum strictly inside \(\TraceCore\) and prove a sharper collapse to it;
    \item do either of the previous two while preserving the same benchmark one-metric package and the same claim boundary.
\end{enumerate}
Another restatement of the full completion-core graph presentation, another restatement of the exact-shell affine nucleus, another same-output deeper-origin relay that merely preserves the same trace core, another reopening of preserved deeper ambient-style relays as if they again defined the live burden, or any move that introduces a second operative metric, enlarges the ordinary matter law, or uses stronger promotion language does not satisfy the present burden.
\end{corollary}

\begin{proof}
Theorem \ref{thm:main} shows that the full completion-core presentation is no longer the minimal benchmark-facing datum at current scope. Therefore a subsequent manuscript must improve on the smaller trace core rather than merely restating the larger completion-core graph presentation or replaying the already-forced nucleus.
\end{proof}

\section{Conclusion}

The positive result of this note is a genuine internal compression below the current completion-core frontier. The previous active theorem already showed that the exact-shell affine nucleus is no longer independently unexplained once the minimal zero-response affine completion core is fixed. The present theorem then takes the second admissible step left open there: it shows that even the full completion-core graph presentation is still larger than necessary.

What survives as the live burden is the smaller completion-state shell-trace core \(\TraceCore\). At this sharpened level, the operative data are only the completion-state body, the affine observer-data and shell-response maps, the hidden body with positive gap operator, the induced symmetry and time-reversal structure, the exact zero-response shell realization in completion-state coordinates, fixed-data solvability, and the same data-kernel coercive control. The affine completion plane and the full completion body are the graph reconstruction of those same trace data, not additional benchmark-facing freedom.

The claim boundary remains fixed throughout. The one operative metric is still exactly \(\CompletedMetric\), the ordinary matter route is unchanged, there is no second operative metric, and the low-energy matter law is not enlarged. This note therefore does not complete a first-principles derivation of GR from substrate microphysics. It does something narrower and honest: it moves the live burden one level further inside the current completion-core frontier. The next benchmark-facing step must now derive or uniquely force that sharper trace core from still more primitive explicit substrate structure, or else collapse it further.

\end{document}
